\begin{textsample}{POS Dim 3 – al – Score 17.00 – al/al\_inf\_17.txt}  \label{ex:f3_pos_218}
\textbf{Crianças} bem \textbf{pequenas} ( 1 ano e 7 \textbf{meses} a 3 anos e 11 \textbf{meses} ) Conforme a teoria do desenvolvimento da pessoa humana de Wallon, entre um e três anos de idade, aproximadamente, a \textbf{criança} vivencia o estágio sensório-motor e projetismo, marcada pela investigação e exploração da realidade exterior, bem como pela aquisição da capacidade simbólica e pelo início da representação.
Ao final do segundo ano, a \textbf{criança} começa a imaginar, sem a necessidade do objeto concreto, ao que chamamos de faz-de-conta.
Os movimentos passam a ser mais controlados, objetivos em relação aos interesses da \textbf{criança}.
As atividades circulares ( \textbf{repetidas} ) servem para reafirmar as novas sensações descobertas.
Essa atividade circular consiste em uma coordenação mútua dos campos \textbf{sensoriais} e motores, possibilitando um ajustamento do \textbf{gesto} a seu efeito, o que refina os progressos de preensão, da \textbf{percepção} e da linguagem, como também permite à \textbf{criança} descobrir as qualidades dos objetos, aguçar as sensibilidades e organizar seus \textbf{gestos} úteis.
Pela organização de suas ações, ela discrimina diferentes movimentos e os efeitos \textbf{sensoriais}, visuais, auditivos e cinestésicos, tornando sua atividade, cada vez mais planejada e organizada voluntariamente ( MAHONEY et al, 2007, p.32 ).
É a fase na qual a \textbf{criança} se coloca de \textbf{pé}, começa a andar e a desenvolver a linguagem oral, momento em que a etapa projetiva do estágio entra em cena, significa que a \textbf{criança}, por meio de suas novas habilidades expressa por meio de \textbf{gestos} e movimentos seu funcionamento mental.
Começa a ver o mundo de um jeito novo.
O movimento e os \textbf{gestos} precedem a palavra, tal como uma \textbf{criança} quando indagada sobre o tamanho do elefante e ela abre os braços para \textbf{demonstrar}.
Graças à função simbólica, isto é, à possibilidade de representar o real, a \textbf{criança} pode elaborar mentalmente o espaço, distribuir os objetos no espaço e no tempo, representá -los e estabelecer signos para as representações. [... ] A função simbólica de desdobramento e de substituição é, portanto, condição para o \textbf{nascimento} do pensamento e da pessoa.
São as atividades de exploração do mundo objetivo que desenvolvem na \textbf{criança} a capacidade de diferenciar objetos dos seres humanos ( MAHONEY et al, 2007, p.32 ).
A \textbf{criança} não é capaz de imaginar sem representar, já é capaz de \textbf{imitar} e simular situações, um prelúdio da representação, que somente se efetiva aproximadamente ao terceiro ano de vida, no estágio do personalismo. • \textbf{Crianças} \textbf{pequenas} ( 4 anos a 5 anos e 11 \textbf{meses} ) O personalismo, que começa aos 3 anos e finaliza aos 6, apresentará mais características neste agrupamento \textbf{etário}.
Nesta fase de construção da consciência de si, através das interações sociais começa o egocentrismo e início da formação da personalidade com a construção da consciência de si, através das interações sociais e relações afetivas.
Nesta idade, há um sentimento forte de posse da \textbf{criança} sobre seus \textbf{brinquedos} e as pessoas que ama, a \textbf{criança} pensa que tudo é apenas dela, inclusive a mãe e o pai e que as coisas acontecem por causa dela.
Por isso, ela precisa compreender que existem pessoas com quem precisará dividir o mundo e as pessoas e esse processo será direcionado nas situações intencionais como um \textbf{dia} de \textbf{partilha} com lanche coletivo, \textbf{partilha} de \textbf{brinquedos}, jogos e \textbf{brincadeiras} com intencionalidades planejadas.
Com isso, a situação muda e o compartilhar fica mais fácil, embora não deva ser forçado ; aos poucos vão se constituindo os vínculos afetivos, surgindo os questionamentos, \textbf{curiosidades} e descobertas, que levam as \textbf{crianças} a desenvolverem representação do seu imaginário.

% matched lemmas: brincadeira, brinquedo, criança, curiosidade, demonstrar, dia, etário, gesto, imitar, mês, nascimento, partilha, pequeno, percepção, pé, repetir, sensorial
\end{textsample}
